This section defines the objects we measure and the explanations we test. Throughout, a
contract pays \$1 if its event occurs. We work with the \emph{side price} $\q\in\{0.01,\dots,0.99\}$:
the price paid for whichever side, YES or NO, a trader buys. Let $\won\in\{0,1\}$ indicate that
the side paid out and let $\pi$ be its true probability.

\subsection{Measuring the bias}\label{sec:framework-measures}

\begin{definition}[Calibration and the favourite--longshot bias]\label{def:flb}
The calibration function is $c(\q)=\E[\won\mid\q]$. Prices exhibit a \emph{favourite--longshot
bias} (FLB) if $c(\q)<\q$ for small $\q$ and $c(\q)>\q$ for large $\q$: longshots are overpriced
and favourites underpriced.
\end{definition}

We summarise the bias in three ways.
\begin{enumerate}
  \item \emph{Returns by price bucket.} The return per dollar of buying a side at $\q$ with fee
  $\fee$ per contract is
  \begin{equation}\label{eq:ret}
    \ret=\frac{\won-\q-\fee}{\q+\fee}.
  \end{equation}
  We report gross ($\fee=0$) and net-of-fee returns.
  \item \emph{Logit slope.} In the regression
  \begin{equation}\label{eq:logit}
    \operatorname{logit}\Prob(\won=1\mid\q)=\alpha+\beta\,\operatorname{logit}\q,
  \end{equation}
  $\beta>1$ indicates an FLB, $\beta<1$ a reverse bias, and $\beta=1,\alpha=0$ calibration.
  \item \emph{Longshot--favourite spread.} The difference in mean return per dollar between
  longshots ($\q\le 10\cents$, \LS) and favourites ($\q\ge 90\cents$, \FAV).
\end{enumerate}

A binary contract has two sides priced $\q$ and $1-\q$ at the same moment. If each observation enters with both sides, the calibration function is symmetric by construction:
$c(1-\q)=1-c(\q)$, so $\alpha=0$ and the bias is carried entirely by $\beta$. If only the YES side
enters, a preference for buying YES ("YES bias") appears as $\alpha<0$. We therefore report
both the side-symmetric and the YES-only versions.

\subsection{Mechanical limits: tick, fees and capital}\label{sec:framework-band}

Some of the bias cannot be traded away even by a risk-neutral, fully informed arbitrageur. Three
frictions matter.
\begin{itemize}
  \item \emph{Tick floor.} Prices live on a one-cent grid. A contract with $\pi<0.01$ cannot trade
  below $1\cents$, so buyers at the floor lose at least $0.01-\pi$ per contract. The loss in the
  lowest bucket is partly mechanical.
  \item \emph{Fees.} A Kalshi taker pays $\fee(\q)=\theta\,\q(1-\q)$ per contract, with
  $\theta=0.07$ times a series multiplier. Makers pay $\theta_m=0.0175$ on some series and
  nothing on others. The fee per contract is the same for both sides of a trade, but per dollar
  it is $\theta(1-\q)$, about 7\% for a longshot and much less for a favourite.
  \item \emph{Capital.} Exploiting overpriced longshots means buying the opposite side at
  $1-\q$ and holding it to resolution, a horizon of $T$ years. This ties up $1-\q$ per contract
  at an opportunity cost $\rho$ per year.
\end{itemize}

\begin{proposition}[No-arbitrage band]\label{prop:band}
Ignore the bid--ask spread. A risk-neutral taker buys the NO side of a longshot priced $\q$ only
if $\q-\pi>\theta\q(1-\q)+\rho T(1-\q)$, and buys the YES side only if
$\pi-\q>\theta\q(1-\q)+\rho T\q$. Prices are therefore consistent with no arbitrage whenever
\begin{equation}\label{eq:band}
  1-\theta(1-\q)-\rho T\,\frac{1-\q}{\q}\;\le\;\frac{\pi}{\q}\;\le\;1+\theta(1-\q)+\rho T .
\end{equation}
\end{proposition}
\begin{proof}
Buying NO at $1-\q$ pays $1-\pi$ in expectation and costs $(1-\q)+\theta\q(1-\q)$ plus
$\rho T(1-\q)$ of capital. Rearranging the condition for a positive expected profit, and dividing
by $\q$, gives the lower bound. The upper bound follows in the same way for buying YES at $\q$,
which ties up $\q$ of capital.
\end{proof}

The lower bound has a feature that matters for longshots. Its capital term scales with
$(1-\q)/\q$, so the band widens sharply as the price falls and as the horizon lengthens.
\begin{itemize}
  \item At $\q=2\cents$, $T=30$ days and $\rho=5\%$, the capital term alone allows the true
  probability to sit about 20\% below the price.
  \item At $\q=50\cents$, the same capital term allows less than 1\%.
\end{itemize}
Mechanical frictions therefore predict a bias concentrated in cheap, long-dated longshots, which
grows with the horizon \citep{pageclemen2013}. Bid--ask spreads widen the band further.

\subsection{Behavioural and informational explanations}\label{sec:framework-theories}

Three classes of explanation predict a bias beyond the band. They make different predictions
about \emph{who} buys longshots.

\paragraph{Preferences: risk-love and probability weighting.} A representative bettor with
convex utility \citep{ali1977} or skewness preference \citep{golec1998} pays a premium for
lottery-like payoffs. Under rank-dependent or cumulative prospect theory
\citep{kahneman1979,tversky1992,jullien2000}, prices reflect decision weights $w(\pi)$ rather than
probabilities. With the linear-in-log-odds weighting function
$\operatorname{logit}w(\pi)=\delta+\gamma\operatorname{logit}\pi$, inverting $\q=w(\pi)$ gives
regression~\eqref{eq:logit} with $\beta=1/\gamma$ and $\alpha=-\delta/\gamma$. A logit slope
above one is thus equivalent to inverse-S weighting ($\gamma<1$), and $\beta$ measures its
strength. \citet{snowbergwolfers2010} argue that misperception of probabilities, rather than
risk-love, best explains racetrack prices.

\paragraph{Information: insiders.} In \citet{shin1991,shin1993}, a market maker who faces a share
$\shinz$ of perfectly informed traders protects himself by shading prices. The shading is
proportionally largest on longshots, because an insider's bet on a longshot is the most
costly to the market maker. For an event with $n$ mutually exclusive outcomes priced
$\q_1,\dots,\q_n$ with sum $S=\sum_j\q_j>1$, the implied probabilities are
\begin{equation}\label{eq:shin}
  \pi_i=\frac{\sqrt{\shinz^2+4(1-\shinz)\,\q_i^2/S}-\shinz}{2(1-\shinz)},
\end{equation}
with $\shinz$ chosen so that $\sum_i\pi_i=1$. The insider share $\shinz$ is a per-event measure
of adverse selection, and it predicts that informed flow, not noise, is on the longshot.

\paragraph{Heterogeneous beliefs.} When risk-neutral traders with dispersed beliefs face budget
constraints \citep{ali1977,gandhi2015}, the most optimistic traders set the price of each outcome.
Since optimism about a longshot moves its price proportionally more, longshots trade above their
mean-belief probability. The bias then rises with belief dispersion and with the share of noise
traders.

\subsection{Testable implications}\label{sec:framework-hypotheses}

The explanations are not mutually exclusive, but they make distinct predictions about which
traders lose on longshots and when.

\begin{hypothesis}[Mechanical]\label{hyp:mech}
If frictions are the whole story, the bias lies inside the band~\eqref{eq:band}. It is
concentrated at the tick floor and grows with horizon, and it shows up in gross but not net
returns.
\end{hypothesis}
\begin{hypothesis}[Preference and retail demand]\label{hyp:pref}
If the bias reflects demand for lottery-like payoffs by unsophisticated traders
\citep{kumar2009,barberis2008}, then:
\begin{itemize}
  \item longshot losses fall on takers, small trades and new or inexperienced accounts;
  \item the bias widens when retail attention is high \citep{barberodean2008,barber2022};
  \item makers and experienced traders earn the other side.
\end{itemize}
\end{hypothesis}
\begin{hypothesis}[Information]\label{hyp:info}
If longshot prices are shaded against insiders, then:
\begin{itemize}
  \item longshot order flow predicts outcomes, at least for large or early trades;
  \item the estimated $\shinz$ is higher where private information is plausible, as in
  politics, announcements and corporate events, than where outcomes are publicly observed, as in
  weather and game scores.
\end{itemize}
\end{hypothesis}
\begin{hypothesis}[Belief dispersion]\label{hyp:beliefs}
If belief heterogeneity drives the bias, it is larger in events with many outcomes, high
turnover and retail-heavy participation, and smaller where a public forecast anchors beliefs.
\end{hypothesis}

\Cref{sec:res-xs} measures the bias against these benchmarks.
\Cref{sec:res-flow,sec:res-wallets} identify who buys longshots. \Cref{sec:res-structural}
compares the explanations directly.
